\chapter{Introduction}
\label{chap:intro}

\section{Background and Problem Context}

The rapid growth of data-intensive applications over the past decade has exposed fundamental limitations in traditional relational database management systems. While RDBMS solutions such as PostgreSQL and MySQL provide robust transactional guarantees and a mature SQL ecosystem, they were designed in an era when workloads were predominantly transactional and datasets fit comfortably on a single machine. Modern applications increasingly demand both transactional consistency for write operations and analytical performance for real-time reads - a combination that forces organizations into maintaining separate systems for each concern, introducing synchronization complexity and operational overhead.

The NoSQL movement emerged as a response to these scalability challenges, offering specialized data models - key-value stores, document databases, graph databases, columnar engines - each optimized for particular access patterns. However, adopting NoSQL typically requires abandoning SQL compatibility, rewriting application code, and accepting weaker consistency models. This creates a difficult trade-off: developers must choose between the familiarity and guarantees of SQL and the performance characteristics of specialized storage engines.

\section{Project Overview}

This diploma project explores a third path - a hybrid architecture that preserves PostgreSQL wire compatibility and a practical SQL subset while leveraging multiple storage technologies internally. The result is Snowdrop, a proof-of-concept database engine that accepts standard PostgreSQL connections and executes SQL queries against a composite backend consisting of a graph database for schema metadata, an embedded key-value store for durable row storage and secondary indexes, and an in-memory columnar cache for analytical read acceleration.

The core hypothesis is that by decomposing the responsibilities of a traditional database engine and assigning each responsibility to the most appropriate storage technology, it is possible to achieve sub-30\,ms p99 query latency on analytical read paths while maintaining ACID guarantees and keeping the PostgreSQL ecosystem usable without changes to clients, drivers, or ORMs.

Snowdrop validates this hypothesis through implementation and benchmarking against PostgreSQL~16 using the TPC-H benchmark at scale factor~1, providing an objective comparison on identical hardware. The implementation is written in Rust, chosen for its zero-cost abstractions, absence of garbage collection pauses, and mature async runtime ecosystem - properties critical for achieving deterministic tail latency in a database engine. In its current state, the system also includes a PostgreSQL-oriented migration workflow with preflight analysis, automatic safe-subset migration, and post-load verification.

\section{Report Structure}

This report is structured as follows. Chapter~1 presents a literature review covering database architecture evolution, hybrid storage approaches, and the theoretical foundations underlying the project's design decisions. Chapter~2 analyzes existing database systems - both traditional RDBMS and modern NewSQL/hybrid solutions - identifying the gaps that Snowdrop aims to address. Chapter~3 describes the data collection strategy, including the TPC-H benchmark dataset and the metrics framework used for performance evaluation. Chapter~4 details the methodology guiding the project's development, from technology selection criteria to benchmark-driven validation. Chapter~5 presents the system architecture through component diagrams, UML diagrams, and the current schema-metadata design in Neo4j. Chapter~6 provides a detailed technology comparison justifying each component choice against evaluated alternatives.
